% =====================================================================
% Section: system implementation and technology stack.
% Sources: requirements.txt, pyproject.toml, docker-compose.yml, src/stages,
% src/utils, experiments/dialstart, be/ (FastAPI service), FE/ (React app).
% =====================================================================

\section{System Implementation and Technology Stack}
\label{sec:system}

\method{} runs in production as a web service behind an inspection interface.
This section describes the technologies used at each layer, from raw transcript to user interface (Table~\ref{tab:stack}), and the engineering choices that the method of \S\ref{sec:method} depends on.

\begin{table*}[t]
\centering
\small
\begin{tabular}{@{}p{0.17\textwidth}p{0.30\textwidth}p{0.46\textwidth}@{}}
\toprule
Layer & Technology & Role in \method \\
\midrule
Pre-processing      & Python 3.11+, rule-based stages            & Apply human corrections with provenance, freeze lines into evidence items with stable IDs, merge lines into speaker turns \\
Topic segmentation  & DialSTART \citep{gao2023dialstart} on ViBERT4news \citep{vibert4news}; PyTorch, Hugging Face Transformers & Default segmenter, Vietnamese checkpoint; TreeSeg \citep{gklezakos2024treeseg} with BGE-M3 embeddings \citep{chen2024bgem3} as alternative \\
Generative LLM      & \texttt{gemma-4-26B-A4B-it} \citep{gemmateam2026gemma4}, OpenAI-compatible gateway, \texttt{openai} SDK & Topic labelling, the four extractors, Verifier and Proposer; strict JSON-schema structured outputs \\
Commitment classifier & OpenAI Decisions API (multiple-choice, calibrated probabilities) & Commitment-act test of the evidence check; status test for progress updates \\
Agent orchestration & LangGraph \citep{langgraph}                & Parent graph and topic subgraph, \texttt{Send} fan-out, list reducers \\
Retrieval tool      & BM25 \citep{robertson2009bm25}, own implementation & \texttt{search\_meeting} over Vietnamese syllables and syllable bigrams \\
Service layer       & FastAPI, Uvicorn, Pydantic v2              & Three report endpoints (tasks, decisions, content) plus cached read-back \\
Result cache        & Redis 7.4 (Docker Compose), \texttt{redis-py} & Two-tier cache keyed by transcript hash and pipeline fingerprint; versions, locks, single-flight \\
User interface      & React 19, TypeScript, Vite, Tailwind CSS 4 & Inspector view: transcript on the left, report on the right, evidence highlighting \\
Testing             & pytest, fakeredis                          & Unit tests for rules, graph routing, actors, deadlines, API and cache \\
\bottomrule
\end{tabular}
\caption{Technology stack of \method, from transcript to user interface. Every component except the two hosted model endpoints runs on a single machine; the generative backbone is open-weight.}
\label{tab:stack}
\end{table*}

% ---------------------------------------------------------------------
\subsection{From transcript to speaker turns}
\label{sec:system:preprocess}

The input is a transcript exported by a speech-to-text system or converted from a Word or plain-text minute: a list of lines, each with an identifier, a speaker and timestamps.
Three deterministic stages prepare it.
Stage~1 overlays human corrections of text or speaker onto the ASR output and records, per field, whether the value comes from ASR or from a reviewer.
Stage~2 freezes each line into an evidence item whose identifier depends only on the meeting and the source line, while the normalised content is hashed separately; re-running the pipeline therefore yields the same citation identifiers, and editing a line changes only its content hash.
Stage~3 merges consecutive lines into a speaker turn only when they share a speaker and lie less than 1.5\,s apart; lines with an unknown speaker or missing timestamps are not merged, because a wrong merge would silently attribute one person's words to another.
Every item of the final report cites these turn identifiers and, through them, the original line identifiers, which the interface uses to highlight evidence.

% ---------------------------------------------------------------------
\subsection{Topic segmentation}
\label{sec:system:segmentation}

The default segmenter is DialSTART \citep{gao2023dialstart}, an unsupervised dialogue topic segmenter that combines a next-sentence-prediction coherence score with the similarity of topic-aware utterance embeddings.
Both of its BERT encoders \citep{devlin2019bert} start from the Vietnamese checkpoint \texttt{vibert4news-base-cased} \citep{vibert4news}, and the model is trained with the DialSTART objective on Vietnamese data.
% TODO: describe the training corpus and settings of checkpoint model_vn3 (not documented in the repo).
For each gap between turns $i$ and $i{+}1$ we compute
$\mathrm{score}(i)=\sigma\big(\mathrm{NSP}(\text{left context}, t_{i+1}) + \cos(\bar e_{\text{left}}, \bar e_{\text{right}})\big)$
over windows of two turns on each side, and derive TextTiling depth scores \citep{hearst1997texttiling}.
Two inference changes make this practical on long meetings: each turn is encoded once and its embedding is averaged per window, instead of re-encoding the window for every gap (equivalent output, about $2\times$ window-size faster); and turns are encoded up to 256 tokens as in training, rather than truncated to 128 characters.
Because the number of boundaries is unknown at inference time, boundaries are selected by smoothing the gap scores (window 3), keeping depths above mean $+$ one standard deviation, requiring at least 4 turns per segment, and dropping segments shorter than a length floor that scales with the meeting (total characters$/15$, clipped to $[1{,}500, 10{,}000]$).
On 12 labelled sessions of the Vietnamese National Assembly, removing the smoothing step lowers boundary F1 from 0.418 to 0.297.
TreeSeg \citep{gklezakos2024treeseg}, a hierarchical segmenter over sentence embeddings, is kept as a configurable alternative and uses BGE-M3 \citep{chen2024bgem3} through \texttt{sentence-transformers} or a hosted embedding API.
% TODO(results): the 0.418 -> 0.297 figure is copied from a comment in
% be/services/segmentation.py; re-run the segmentation eval before citing it.

% ---------------------------------------------------------------------
\subsection{LLM access and structured outputs}
\label{sec:system:llm}

All generative roles call one open-weight model, \texttt{gemma-4-26B-A4B-it} \citep{gemmateam2026gemma4}, a mixture-of-experts model served behind an OpenAI-compatible endpoint and accessed with the \texttt{openai} Python SDK.
Every call goes through a single adapter interface, \texttt{generate\_json(system, user, schema)}, which requests strict JSON-schema structured outputs at temperature~0; a looser ``valid JSON'' mode is available for servers without schema-constrained decoding.
Because every node only depends on this interface, the Verifier can be pointed at a different provider (e.g.\ a hosted DeepSeek model) by configuration, with its own concurrency limit and without code changes.
The adapter retries transient failures (timeouts, dropped connections, HTTP 5xx) with exponential backoff, up to 4 attempts, 6 for connection errors and 8 for rate limits, with a 60\,s request timeout.
A context variable labels each call with the agent and segment that made it, so retries and latency can be attributed even though many calls run concurrently.
Since all segments run in parallel, the number of simultaneous requests could exceed what the gateway tolerates; a process-wide semaphore (\S\ref{sec:system:orchestration}) bounds it.

The commitment classifier of \S\ref{sec:method:evidence} uses the OpenAI Decisions API, which answers a multiple-choice question about an input text with a probability for each choice.
Results are cached by input text, because several candidates of one topic often point at the same commitment turn, and a 10\,s timeout falls back to the cue lexicon.
On the development turns its median latency is 0.30\,s (p95 0.41\,s), small compared with a generative call.

% ---------------------------------------------------------------------
\subsection{Agent orchestration}
\label{sec:system:orchestration}

The pipeline is a LangGraph \citep{langgraph} \texttt{StateGraph}.
The parent graph has three nodes (Planner, topic subgraph, Finalize); the topic subgraph is compiled separately with its own input and output schemas, so it receives one segment and returns only its result lists.
Fan-out happens twice with \texttt{Send}: once per segment and once per UNCERTAIN candidate.
Results from concurrent branches are combined with list-concatenation reducers and re-sorted in Finalize.
The graph has no human-in-the-loop interrupt and runs from start to end in one invocation, so it needs no checkpointer, and one compiled graph serves all meetings.
Graph parallelism and LLM load are controlled separately: LangGraph runs up to 16 tasks per superstep, while a bounded semaphore wrapped around every LLM adapter (labeller, extractors, Verifier) caps concurrent LLM calls at 8 by default, below the client's connection pool of 10.
Each extractor produces exactly one report section, so a caller can recompute one section on top of stored topic labels without touching the others.
Table~\ref{tab:hparams} lists the pipeline hyperparameters.

\begin{table}[t]
\centering
\small
\begin{tabular}{@{}llr@{}}
\toprule
Symbol & Meaning & Value \\
\midrule
$A$ & attempts per extractor and topic & 2 \\
$T$ & Verifier tool calls per round & 3 \\
$R$ & max.\ consensus rounds & 3 \\
$k$ & turns returned by \texttt{search\_meeting} & 5 \\
$k_1, b$ & BM25 parameters & 1.2, 0.75 \\
-- & commitment bands of $p_{\text{commit}}$ & 0.2 / 0.8 \\
-- & context turns for the classifier & 2 \\
-- & amendment retention (min.) & 0.5 \\
-- & actor clear: score / margin & 0.6 / 0.2 \\
-- & duplicate: Jaccard / overlap & 0.45 / 0.7 \\
-- & graph tasks per superstep & 16 \\
-- & concurrent LLM calls (default) & 8 \\
\bottomrule
\end{tabular}
\caption{Hyperparameters of \method{} (defaults of the released configuration). With these values a single UNCERTAIN candidate costs at most $R(T{+}3)=18$ LLM calls.}
\label{tab:hparams}
\end{table}

% ---------------------------------------------------------------------
\subsection{Service layer and caching}
\label{sec:system:service}

A FastAPI application (served by Uvicorn, request and response models in Pydantic v2) exposes the report as three documents: \emph{tasks} (new action items and progress of agenda tasks), \emph{decisions}, and \emph{content} (topics with discussion points).
Each is a \texttt{POST} endpoint that returns its document in the same request; a \texttt{GET} endpoint returns a document that has already been computed.
Every document item carries the identifiers of the transcript lines it cites, together with an audit block built from the verification records.

Results are cached in two tiers under one key, a hash of the normalised transcript together with a \emph{pipeline fingerprint}: a hash of the source code, configuration files and model or endpoint names (never API keys), so editing a prompt invalidates old results even before it is committed.
The \emph{base} tier holds speaker turns, segments and topic labels and is shared by the three documents; each document is a separate entry that is computed only when requested, on top of the base.
The default backend is Redis~7.4, run with Docker Compose and append-only persistence; each analysis is a Redis hash with a time-to-live.
Version counters let the interface discard a response older than the one it shows, locks keep read-modify-write sections consistent, and a single-flight guard makes a duplicate request wait for the running computation instead of starting another, which also holds across several API workers.
Disk, in-memory (LRU) and disabled backends exist for development.

% ---------------------------------------------------------------------
\subsection{Inspection interface}
\label{sec:system:ui}

The web client is written in React~19 and TypeScript, built with Vite and styled with Tailwind CSS~4.
It follows an inspector layout: the transcript stays pinned on the left and the three documents render on the right.
Clicking any action item, decision or discussion point highlights the exact transcript lines it cites, so a reader can check each claim against its source; this is the user-facing side of goal G3 (\S\ref{sec:method}).
